\documentclass[11pt,a4paper]{article}
\usepackage[utf8]{inputenc}
\usepackage[spanish]{babel}
\usepackage[margin=2.5cm]{geometry}
\usepackage{amsmath,amssymb}
\usepackage[version=4]{mhchem} % Para formulas quimicas
\usepackage{array}
\usepackage{booktabs}
\usepackage{xcolor}
\usepackage{tcolorbox}
\usepackage{enumitem}
\usepackage{fancyhdr}
\usepackage{titlesec}

% Colores personalizados
\definecolor{azuloscuro}{RGB}{0,51,102}
\definecolor{verde}{RGB}{0,128,0}
\definecolor{rojo}{RGB}{180,0,0}
\definecolor{amarillo}{RGB}{255,200,0}
\definecolor{grisclaro}{RGB}{245,245,245}

% Configuracion de secciones
\titleformat{\section}{\Large\bfseries\color{azuloscuro}}{\thesection}{1em}{}
\titleformat{\subsection}{\large\bfseries\color{azuloscuro!80}}{\thesubsection}{1em}{}

% Cajas para resultados
\newtcolorbox{resultado}[1][]{
    colback=grisclaro,
    colframe=azuloscuro,
    fonttitle=\bfseries,
    title=#1,
    boxrule=0.5pt,
    arc=3pt
}

\newtcolorbox{reaccionbox}{
    colback=white,
    colframe=verde!70!black,
    boxrule=0.5pt,
    arc=2pt,
    left=5pt,
    right=5pt
}

% Encabezado y pie de pagina
\pagestyle{fancy}
\fancyhf{}
\fancyhead[L]{\small Lab 1Q.A. 2025-10}
\fancyhead[R]{\small Identificacion de Aniones y Cationes}
\fancyfoot[C]{\thepage}
\renewcommand{\headrulewidth}{0.4pt}

\begin{document}

% ===========================================
% PORTADA
% ===========================================
\begin{center}
    {\LARGE\bfseries\color{azuloscuro} Laboratorio de Quimica Analitica}\\[0.5cm]
    {\Large Lab 1Q.A. 2025-10 (USS)}\\[1cm]
    {\Huge\bfseries Identificacion de Aniones y Cationes}\\[1.5cm]
    \rule{\textwidth}{1pt}
\end{center}

\vspace{1cm}

% ===========================================
% MATERIALES
% ===========================================
\section{Materiales Necesarios}

\begin{itemize}[leftmargin=*]
    \item \ce{HNO3} 6M (2 unidades)
    \item \ce{HCl} 6M (2 unidades)
    \item 11 tubos de ensayo
    \item Trabajar en 4 estaciones, 8 grupos
\end{itemize}

% ===========================================
% PREPARACION
% ===========================================
\section{Preparacion de Muestras}

\subsection{Aniones}
\textbf{4 grupos $\rightarrow$ ANIONES:} Son 11 aniones. Sacar una punta de espatula, disolver en $\sim$40g de agua (aproximadamente 2 dedos de altura en el tubo).

\vspace{0.3cm}
\textbf{Aniones a identificar:}
\begin{center}
\ce{PO4^{3-}}, \ce{SO4^{2-}}, \ce{SO3^{2-}}, \ce{I^-}, \ce{Br^-}, \ce{Cl^-}, \ce{NO2^-}, \ce{NO3^-}, \ce{CO3^{2-}}, \ce{C2O4^{2-}}, \ce{CrO4^{2-}}
\end{center}

\textbf{Nota:} Disolucion dificil con \ce{C2O4^{2-}} (oxalato)

\subsection{Cationes (9 tubos)}

\begin{center}
\begin{tabular}{clcc}
\toprule
\textbf{N°} & \textbf{Compuesto} & \textbf{Cation} & \textbf{Color} \\
\midrule
1 & \ce{BaCl2} & \ce{Ba^{2+}} & -- \\
2 & \ce{AgNO3} & \ce{Ag^+} & -- \\
3 & \ce{Ca(NO3)2} & \ce{Ca^{2+}} & -- \\
4 & \ce{FeCl3 * 6H2O} & \ce{Fe^{3+}} & Amarillo \\
5 & \ce{NiCl2 * 6H2O} & \ce{Ni^{2+}} & Verde \\
6 & \ce{AlCl3} & \ce{Al^{3+}} & -- \\
7 & \ce{ZnCl2} & \ce{Zn^{2+}} & -- \\
8 & \ce{CoCl2 * 6H2O} & \ce{Co^{2+}} & Rojo \\
9 & \ce{Cu(NO3)2 * 3H2O} & \ce{Cu^{2+}} & Celeste \\
\bottomrule
\end{tabular}
\end{center}

\subsection{Notas Importantes}
\begin{itemize}[leftmargin=*]
    \item Fijarse en colores de las soluciones
    \item Limpiar espatula entre cada muestra
    \item \ce{NH3} = \ce{NH4OH} (en solucion acuosa)
    \item Lavar vasos y tubos despues de cada ensayo
    \item Acetato de etilo = alcohol etilico
\end{itemize}

\newpage

% ===========================================
% ENSAYOS DE ANIONES
% ===========================================
\section{Ensayos Especificos -- ANIONES}

% --- 1. Fosfato ---
\subsection{1. Fosfato (\ce{PO4^{3-}})}

\textbf{Reactivo:} Molibdato de amonio \ce{(NH4)6Mo7O24 * 4H2O}

\textbf{Procedimiento:}
\begin{enumerate}[leftmargin=*]
    \item Agregar 5 mL de molibdato de amonio 0.5 mol/L
    \item Agregar 4 mL de muestra
    \item Agregar 1 mL de \ce{HNO3} 6M (medio acido)
\end{enumerate}

\begin{reaccionbox}
\textbf{Reaccion:}
\begin{equation*}
\ce{(NH4)6Mo7O24 + PO4^{3-} + H^+ -> (NH4)3[PMo12O40] + H2O + NH4^+}
\end{equation*}
\end{reaccionbox}

\begin{resultado}[Resultado positivo]
Precipitado \textbf{\textcolor{amarillo!80!black}{AMARILLO}}
\end{resultado}

% --- 2. Sulfato ---
\subsection{2. Sulfato (\ce{SO4^{2-}})}

\textbf{Reactivo:} Cloruro de bario \ce{BaCl2}

\textbf{Procedimiento:}
\begin{enumerate}[leftmargin=*]
    \item Agregar gotas de \ce{BaCl2}
    \item Agregar \ce{H^+} (\ce{HCl} 6M, 3 gotas)
\end{enumerate}

\begin{reaccionbox}
\textbf{Reaccion:}
\begin{equation*}
\ce{SO4^{2-} + Ba^{2+} -> BaSO4 v}
\end{equation*}
\end{reaccionbox}

\begin{resultado}[Resultado positivo]
Precipitado \textbf{BLANCO}
\end{resultado}

% --- 3. Sulfito ---
\subsection{3. Sulfito (\ce{SO3^{2-}})}

\textbf{Reactivos:}
\begin{itemize}[leftmargin=*]
    \item 5g \ce{ZnSO4} saturado
    \item 5g \ce{K4[Fe(CN)6]} (Ferrocianuro de potasio)
    \item 2g \ce{Na2[Fe(CN)5NO]} (Nitroprusiato de sodio)
\end{itemize}

\textbf{Procedimiento:}
\begin{enumerate}[leftmargin=*]
    \item Mezclar \ce{ZnSO4} con ferrocianuro $\rightarrow$ \ce{Zn2[Fe(CN)6] v}
    \item Agregar nitroprusiato de sodio
    \item Mezclar con la muestra
\end{enumerate}

\begin{resultado}[Resultado positivo]
Precipitado \textbf{\textcolor{rojo}{ROJO}} (contiene \ce{SO3^{2-}})
\end{resultado}

% --- 4. Yoduro ---
\subsection{4. Yoduro (\ce{I^-})}

\textbf{Reactivos:} \ce{KI}, \ce{NO2^-}, \ce{CH3COOH}, \ce{CHCl3} (cloroformo)

\textbf{Procedimiento:}
\begin{enumerate}[leftmargin=*]
    \item \textbf{Trabajar bajo campana de extraccion}
    \item Agregar 2-3g de \ce{CH3COOH}
    \item Se libera \ce{I2} (yodo)
    \item Extraer con \ce{CHCl3} ($d = 1.49$ g/mL)
\end{enumerate}

\begin{reaccionbox}
\textbf{Reaccion:}
\begin{equation*}
\ce{KI_{(ac)} + NO2^- + H^+ -> I2 + NO ^  + K^+ + H2O}
\end{equation*}
\end{reaccionbox}

\begin{resultado}[Resultado positivo]
Color \textbf{CAFE} (gas violeta). \ce{I2} en \ce{CHCl3} $\rightarrow$ color \textbf{\textcolor{violet}{LILA}}
\end{resultado}

% --- 5. Bromuro ---
\subsection{5. Bromuro (\ce{Br^-})}

\textbf{Reactivos:} \ce{MnO4^-} (permanganato), \ce{H2SO4} 3M, \ce{CHCl3}

\textbf{Procedimiento:}
\begin{enumerate}[leftmargin=*]
    \item \textbf{Trabajar bajo campana} (gases)
    \item Agregar 4g de \ce{H2SO4} 3M
    \item Extraer \ce{Br2} con \ce{CHCl3}
\end{enumerate}

\begin{reaccionbox}
\textbf{Reaccion:}
\begin{equation*}
\ce{2Br^- + MnO4^- + 8H^+ -> Br2 ^ + Mn^{2+} + 4H2O}
\end{equation*}
\end{reaccionbox}

\begin{resultado}[Resultado positivo]
Color \textbf{\textcolor{rojo}{ROJIZO}}. \ce{Br2} en \ce{CHCl3} $\rightarrow$ \textbf{\textcolor{rojo}{ROJIZO}}
\end{resultado}

% --- 6. Cloruro ---
\subsection{6. Cloruro (\ce{Cl^-})}

\textbf{Reactivo:} \ce{AgNO3} (Nitrato de plata)

\begin{reaccionbox}
\textbf{Reaccion:}
\begin{equation*}
\ce{Cl^- + Ag^+ -> AgCl v}
\end{equation*}
\end{reaccionbox}

\begin{resultado}[Resultado positivo]
Precipitado \textbf{BLANCO}
\end{resultado}

% --- 7. Nitrito ---
\subsection{7. Nitrito (\ce{NO2^-})}

\textbf{Reactivos:} \ce{HCl} 6 mol/L, urea \ce{CO(NH2)2}

\textbf{Procedimiento:}
\begin{enumerate}[leftmargin=*]
    \item Agregar 10g de \ce{HCl} 6 mol/L
    \item Agregar punta de espatula de urea
\end{enumerate}

\begin{reaccionbox}
\textbf{Reaccion:}
\begin{equation*}
\ce{CO(NH2)2 + 2HNO2 -> 2N2 ^ + CO2 ^ + 3H2O}
\end{equation*}
\end{reaccionbox}

\begin{resultado}[Resultado positivo]
Desprendimiento de gas \ce{N2} (burbujas)
\end{resultado}

% --- 8. Nitrato ---
\subsection{8. Nitrato (\ce{NO3^-})}

\textbf{Reactivos:} \ce{CH3COOH} 6 mol/L, Zn en polvo, \ce{HCl} 6 mol/L, urea

\textbf{Procedimiento:}
\begin{enumerate}[leftmargin=*]
    \item Agregar 10g de \ce{CH3COOH} 6 mol/L
    \item Agregar Zn en polvo
    \item Esperar 3-4 minutos
    \item Agregar 10g \ce{HCl} 6 mol/L + punta de espatula de urea
\end{enumerate}

\begin{resultado}[Resultado positivo]
Desprendimiento de gas \ce{N2} (burbujas)
\end{resultado}

% --- 9. Carbonato ---
\subsection{9. Carbonato (\ce{CO3^{2-}})}

\textbf{Reactivo:} \ce{CH3COOH} 6 mol/L (acido acetico)

\textbf{Procedimiento:} Agregar 3-4g de \ce{CH3COOH} 6 mol/L

\begin{reaccionbox}
\textbf{Reaccion:}
\begin{equation*}
\ce{CO3^{2-} + 2H^+ -> CO2 ^ + H2O}
\end{equation*}
\end{reaccionbox}

\begin{resultado}[Resultado positivo]
Desprendimiento de gas \ce{CO2} (burbujas/efervescencia)
\end{resultado}

% --- 10. Oxalato ---
\subsection{10. Oxalato (\ce{C2O4^{2-}})}

\textbf{Reactivos:} \ce{KMnO4} (permanganato de potasio), \ce{H2SO4}

\textbf{Procedimiento:}
\begin{enumerate}[leftmargin=*]
    \item Agregar gotas de \ce{C2O4^{2-}}
    \item Agregar 10g de \ce{KMnO4}
    \item Agregar 2-3g de \ce{H2SO4}
    \item \textbf{NO} agregar mas de 2 mL
\end{enumerate}

\begin{reaccionbox}
\textbf{Reaccion:}
\begin{equation*}
\ce{5C2O4^{2-} + 2MnO4^- + 16H^+ -> 10CO2 ^ + 2Mn^{2+} + 8H2O}
\end{equation*}
\end{reaccionbox}

\begin{resultado}[Resultado positivo]
Cambio de color \textbf{\textcolor{violet}{VIOLETA}} $\rightarrow$ \textbf{INCOLORO} (decoloracion del permanganato)
\end{resultado}

% --- 11. Cromato ---
\subsection{11. Cromato (\ce{CrO4^{2-}})}

\textbf{Reactivos:} \ce{H2SO4}, acetato de etilo, \ce{H2O2}

\textbf{Procedimiento:}
\begin{enumerate}[leftmargin=*]
    \item Agregar 2-3g de \ce{H2SO4} 6g/L
    \item Agregar 1 mL de acetato de etilo
    \item Agregar 2-3g de \ce{H2O2}
    \item El alcohol amilico ayuda a estabilizar
    \item \textbf{Nota:} \ce{CrO5} es poco estable
\end{enumerate}

\begin{reaccionbox}
\textbf{Reaccion:}
\begin{equation*}
\ce{CrO4^{2-} + 2H^+ + H2O2 -> CrO5 + 3H2O}
\end{equation*}
\end{reaccionbox}

\begin{resultado}[Resultado positivo]
Capa organica color \textbf{\textcolor{blue}{AZUL}}
\end{resultado}

\newpage

% ===========================================
% ENSAYOS DE CATIONES
% ===========================================
\section{Ensayos Especificos -- CATIONES}

% --- 1. Plata ---
\subsection{1. Plata (\ce{Ag^+})}

\textbf{Fuente:} \ce{AgNO3} (Nitrato de plata)

\textbf{Reactivo:} \ce{KI} (Yoduro de potasio) 0.5 mol/L

\textbf{Procedimiento:}
\begin{enumerate}[leftmargin=*]
    \item Agregar 10g de muestra
    \item Agregar \ce{KI} gota a gota
\end{enumerate}

\begin{reaccionbox}
\textbf{Reaccion:}
\begin{equation*}
\ce{Ag^+ + I^- -> AgI v}
\end{equation*}
\end{reaccionbox}

\begin{resultado}[Resultado positivo]
Precipitado \textbf{\textcolor{amarillo!80!black}{AMARILLO}}
\end{resultado}

% --- 2. Bario ---
\subsection{2. Bario (\ce{Ba^{2+}})}

\textbf{Reactivo:} \ce{K2CrO4} (Cromato de potasio)

\textbf{Procedimiento:}
\begin{enumerate}[leftmargin=*]
    \item Agregar 3g de muestra
    \item Agregar 10g de \ce{K2CrO4} (en vaso de precipitado)
\end{enumerate}

\begin{reaccionbox}
\textbf{Reaccion:}
\begin{equation*}
\ce{Ba^{2+} + CrO4^{2-} -> BaCrO4 v}
\end{equation*}
\end{reaccionbox}

\begin{resultado}[Resultado positivo]
Precipitado \textbf{\textcolor{amarillo!80!black}{AMARILLO}}
\end{resultado}

% --- 3. Calcio ---
\subsection{3. Calcio (\ce{Ca^{2+}})}

\textbf{Reactivo:} \ce{NaOH} 6 mol/L

\textbf{Procedimiento:}
\begin{enumerate}[leftmargin=*]
    \item Agregar 10g de muestra
    \item Agregar \ce{NaOH} gota a gota
\end{enumerate}

\begin{reaccionbox}
\textbf{Reaccion:}
\begin{equation*}
\ce{Ca^{2+} + 2OH^- -> Ca(OH)2 v}
\end{equation*}
\end{reaccionbox}

\begin{resultado}[Resultado positivo]
Precipitado \textbf{BLANCO GELATINOSO}
\end{resultado}

% --- 4. Hierro III ---
\subsection{4. Hierro III (\ce{Fe^{3+}})}

\textbf{Reactivo:} \ce{KSCN} (Tiocianato de potasio) 0.1 mol/L

\textbf{Procedimiento:}
\begin{enumerate}[leftmargin=*]
    \item Agregar 10g de muestra
    \item Agregar 5-10g de \ce{KSCN}
    \item \textbf{IMPORTANTE: NO usar \ce{HCl}}
\end{enumerate}

\begin{reaccionbox}
\textbf{Reaccion:}
\begin{equation*}
\ce{Fe^{3+} + SCN^- -> [FeSCN]^{2+}}
\end{equation*}
\end{reaccionbox}

\begin{resultado}[Resultado positivo]
Solucion color \textbf{\textcolor{rojo}{ROJO}} intenso
\end{resultado}

% --- 5. Aluminio ---
\subsection{5. Aluminio (\ce{Al^{3+}})}

\textbf{Reactivo:} \ce{NH4OH} (Hidroxido de amonio)

\textbf{Procedimiento:}
\begin{enumerate}[leftmargin=*]
    \item Agregar 10g de muestra
    \item Agregar \ce{NH4OH}
\end{enumerate}

\begin{reaccionbox}
\textbf{Reaccion:}
\begin{equation*}
\ce{Al^{3+} + 3NH4OH -> Al(OH)3 v + 3NH4^+}
\end{equation*}
\end{reaccionbox}

\begin{resultado}[Resultado positivo]
Precipitado \textbf{BLANCO GELATINOSO}
\end{resultado}

% --- 6. Zinc ---
\subsection{6. Zinc (\ce{Zn^{2+}})}

\textbf{Reactivo:} \ce{NaOH} (Hidroxido de sodio)

\textbf{Procedimiento:} Agregar 30 gotas de \ce{NaOH}

\begin{reaccionbox}
\textbf{Reaccion:}
\begin{equation*}
\ce{Zn^{2+} + 2OH^- -> Zn(OH)2 v}
\end{equation*}
\end{reaccionbox}

\begin{resultado}[Resultado positivo]
Precipitado \textbf{BLANCO GELATINOSO}
\end{resultado}

% --- 7. Cobre ---
\subsection{7. Cobre II (\ce{Cu^{2+}})}

\textbf{Reactivo:} \ce{NH3} (Amoniaco)

\textbf{Procedimiento:}
\begin{enumerate}[leftmargin=*]
    \item Agregar 10g de muestra
    \item Agregar \ce{NH3} gota a gota
\end{enumerate}

\begin{reaccionbox}
\textbf{Reaccion:}
\begin{equation*}
\ce{Cu^{2+} + 4NH3 -> [Cu(NH3)4]^{2+}}
\end{equation*}
\end{reaccionbox}

\begin{resultado}[Resultado positivo]
Solucion color \textbf{\textcolor{blue}{AZUL REY}} (complejo tetraaminocobre)
\end{resultado}

% --- 8. Niquel ---
\subsection{8. Niquel II (\ce{Ni^{2+}})}

\textbf{Reactivo:} DMG (Dimetilglioxima) en \ce{NH3}

\textbf{Procedimiento:}
\begin{enumerate}[leftmargin=*]
    \item Agregar 10g de muestra
    \item Agregar DMG en medio de \ce{NH3}
\end{enumerate}

\begin{reaccionbox}
\textbf{Reaccion:}
\begin{equation*}
\ce{Ni^{2+} + 2DMG -> Ni(DMG)2 v}
\end{equation*}
\end{reaccionbox}

\begin{resultado}[Resultado positivo]
Precipitado \textbf{\textcolor{red!70!black}{ROJO (FUCSIA)}}
\end{resultado}

% --- 9. Cobalto ---
\subsection{9. Cobalto II (\ce{Co^{2+}})}

\textbf{Reactivos:} \ce{KSCN} (Tiocianato de potasio), \ce{NH4OH}

\textbf{Procedimiento:}
\begin{enumerate}[leftmargin=*]
    \item Agregar 10g de muestra
    \item Agregar \ce{KSCN} (solido)
    \item Agregar \ce{NH4OH}
    \item \textbf{Nota:} \ce{NH3} = \ce{NH4OH}
\end{enumerate}

\begin{reaccionbox}
\textbf{Reaccion:}
\begin{equation*}
\ce{Co^{2+} + 4SCN^- -> [Co(SCN)4]^{2-}}
\end{equation*}
\end{reaccionbox}

\begin{resultado}[Resultado positivo]
Solucion color \textbf{\textcolor{blue}{AZUL}}
\end{resultado}

\vfill

\begin{center}
\rule{0.5\textwidth}{0.5pt}\\[0.3cm]
\textit{Documento digitalizado a partir de notas de laboratorio}\\
\textit{Lab 1Q.A. 2025-10 -- Universidad San Sebastian}
\end{center}

\end{document}
